\section{Implications and Scope}\label{sec:implications}

\subsection{Scope}\label{sec:scope}

The results hold under seven conditions.
\begin{enumerate}[label=(\roman*)]
\item \emph{Security dependence.} The holder relies on its patron for protection it cannot replace on the relevant horizon ($S_k>0$).
\item \emph{A partly conditional guarantee.} Some of that protection can be withheld after exit at a cost the patron can stake, and making all of it conditional carries an assurance cost ($\kappa S_k+d_p>0$; \cref{prop:conditionality}).
\item \emph{A rival that bids.} A second great power can credibly promise something for the chokepoint ($\omega>0$).
\item \emph{A nested chokepoint.} The chokepoint sits partly inside the patron's jurisdiction, so flow jurisdiction is cheap ($\theta>0$), but the holder's licence still decides whether exit happens ($\theta<1$).
\item \emph{Complete information} about conditionality.
\item \emph{Unitary great powers.}
\item \emph{One rival}, with holders bargaining one at a time.
\end{enumerate}
The Netherlands meets the first four (\cref{sec:background}), and \cref{sec:decide} relaxes the unitary holder. On our reading Japan and South Korea meet them too. Both are treaty allies; Beijing threatened retaliation when Washington pressed Tokyo to restrict its toolmakers' sales and servicing in China \citep{taipeitimes2024retaliation}, and Korea's fabs in China run on US-origin tools whose validated end-user status BIS withdrew in 2025 \citep{bis2025veu}. Taiwan fails the third condition: China threatens its security and does not bid for its alignment, so Washington's 2025 proposal that half the chips it uses be made in the United States, which Taipei rejected \citep{younger2025fiftyfifty, focustaiwan2025fiftyfifty}, is a demand for control without a bidder. The Gulf partners that traded security cooperation and investment for chip access \citep{whitehouse2025uae, commerce2025gulfchips, bis2026uae} have no treaty guarantee and fail the first. Kyrgyzstan, which used a Russian aid package to raise the rent for the Manas base \citep{cooley2012great}, fits only the bidding benchmark of \cref{lem:bidding}. \Cref{sec:cases} compares the cases. Because conditionality is a choice (\cref{sec:conditionality}), a less transactional administration would move Europe to another point in the parameter space and leave it within the scope.

\subsection{Implications for Europe}

Each implication holds in a region of the parameter space. On our reading Europe has a nonpositive premium and slack to spend: The Hague has received revocable statuses, which we do not code as transfers (\cref{sec:coding}), and Washington has tied favourable treatment and AI chip exports to export-control alignment, the linkage \cref{prop:counter}(b) leads us to expect against an ally with slack (\cref{sec:counter}). That reading rests on the jurisdictional channel. How much of the slack comes from the conditional part of the guarantee, the record leaves untested (\cref{sec:evidence}).

\paragraph{Reduce the conditional part of dependence.} At The Hague in June 2025 NATO allies committed to spend 5 percent of GDP a year by 2035, at least 3.5 percent on core defence \citep{nato2025hague}. The target counts spending. In the model a euro raises Europe's premium only insofar as it replaces something Washington could withhold after exit, which points to substitutes for posture and enablers (intelligence, munitions, sustainment, forces that can be moved one ally at a time) more than to duplicates of the statutory core (\cref{sec:security}). Below zero, each unit of premium gained removes a unit of the slack Washington can spend (\cref{cor:rationing}), though Europe still receives no transfer. Washington can compete the gain away. Where assurance is cheap at the margin, it answers a higher premium with a more conditional guarantee and pays nothing (\cref{prop:conditionality}(c)(i)); Europe gains only where assurance is already costly, where Washington is already buying slack to spend, or where conditionality is at its ceiling. If public US statements tying protection to European conduct multiply as European capability grows, Europe is in the first regime.

\paragraph{Reduce the nesting.} Moving US-origin content and US-person servicing out of US jurisdiction lowers $\theta$ and raises China's credible offer, but only while China's commitment cap does not bind (\cref{eq:omega}). Jurisdiction is a choice variable \citep{kalyanpur2019mobilizing}, but changing it is slow and visible. On the servicing side, de-nesting lowers the share $h$ that Washington can compel, which is what gives hold-up its value (\cref{sec:holdup}). We do not know whether the cap binds; Chinese offers limited by what Beijing can credibly deliver, and not by US reach, would show that it does, and then de-nesting buys nothing.

\paragraph{Prepare an interruption only as a standing rule.} Interrupting servicing raises Europe's transfer only if a rule fixed before offers are made decides it and keeps it in force for the dispute; decided after a rejected offer, it is never used (\cref{cor:interrupt}). A response measure under the Anti-Coercion Instrument \citep{eu2023aci}, once adopted through the stages of \cref{sec:council}, could be such a rule. It helps only for servicing Washington cannot compel, which for US-sited fabs is a small share on our reading of ASML's US structure (\cref{sec:holdup}).

\paragraph{Expect a preventive demand.} Autonomy investments that Europe cannot commit to forgo strengthen Washington's motive to demand control while its threats still bite, and with investment Washington can buy control with no bite at all (\cref{cor:control}(a)--(b), \cref{cor:windows}). Whether the window opens depends on the net drift of the premium, which rearmament raises and Chinese indigenization lowers (\cref{sec:windows}). Sovereignty costs fixed in advance (ASML's preference-share foundation, Vifo screening, the Foreign Subsidies Regulation) deter the demand (\cref{sec:controlsetup}). Demands for relocation, board seats or a licensing veto that follow visible European investment would show that Europe is in the window (\cref{sec:predictions}).

\paragraph{Bind, or keep the key.} Binding against exit pays only when it pre-empts a demand that would otherwise come, at lower cost, and it pre-empts only if Washington values control little for its own sake; a partial binding can invite a demand by raising the slack Washington's threats draw on (\cref{rem:bind}). The EU has moved partway. The 2023 Economic Security Strategy flagged advanced semiconductors for possible new export tools \citep{eu2023strategy}, the 2024 White Paper on Export Controls, COM(2024) 25, proposed uniform EU controls and a senior-level coordination forum \citep{ec2024whitepaper}, and Delegated Regulation 2025/2003 put the most advanced immersion tools on the EU list (\cref{sec:bind}). On our reading these raise $\Delta_b$ without approaching $\Delta_b=\infty$, and whether that weakened the US motive turns on $v_c$. A shift in US demands from alignment of controls toward ownership or a licensing veto would indicate a high $v_c$.

\paragraph{Chinese counter-threats.} The Critical Raw Materials Act sets 2030 benchmarks of 10 percent EU extraction, 40 percent processing and 25 percent recycling of strategic raw materials, with at most 65 percent from any one third country \citep{eu2024crma}. In the model, removing China's counter-threats weakly raises Washington's payoff and leaves Europe's loss at $B(J^c)$ when Washington has demands China cannot reach; otherwise Europe's loss can rise or fall (\cref{prop:counter}(a)). It also weakens the Chinese threat that protects European ownership against a transfer (\cref{cor:control}(d)). Europe's gain from de-risking therefore lies mostly in disputes over its own measures, such as the Nuctech investigation (\cref{sec:counter}) and Nexperia (\cref{sec:record}). The Blocking Statute \citep{eu1996blocking} lists no US export-control law \citep{ec2018blocking}, so its penalty is zero today (\cref{rem:firm}); listing such laws would matter only once the premium is positive and Europe licenses exit. We treat China's coercive reach and the credibility of its promises as separate. If Europe discounts the promises of a China able to cut its supplies, de-risking also raises $\lambda_C$ and with it $\omega$.

\subsection{Lessons for Washington}

Conditionality is costly. A patron with no further demands makes its guarantee just conditional enough to hold the premium at zero, and more conditionality pays only if the slack is spent, while other allies pay through weaker assurance (\cref{prop:conditionality}). Demands across domains (tariffs, the spending target, procurement, servicing in China) share one budget, so a stronger Chinese bid forces Washington to drop its lowest-value demands first (\cref{cor:rationing}). Assurance raises what coercion yields. Conditional parts are measured from what a loyal ally receives, so untargeted drawdowns and rationing of loyal allies lower $d_p$ and $d_{ai}$ and with them Europe's slack (\cref{sec:conditionality}), and a targeted cut that is carried out spends the instrument (\cref{prop:counter}(b)).

\subsection{Limitations}

Complete information about conditionality is the strongest assumption. If $\kappa$ were private, Washington would gain from overstating it, European beliefs would set $Q$, and documents such as the 2025 National Security Strategy would be partly cheap talk. Refusals and executed counter-threats could then come from misjudgment, the reading of Nexperia that lies outside the model (\cref{sec:record}). The assurance cost $H$ is a reduced form, and who pays for China's bid depends on its shape (\cref{prop:conditionality}(c)); if it reflects Washington's reputation with Tokyo and Seoul, conditionality toward Europe depends on what other holders observe. The great powers are unitary, although Congress can set stakes the executive would not, as the MATCH Act's 150-day deadline for allies would (\cref{sec:background}), and member-specific exposure enters only the Council remark (\cref{rem:council}). With one rival and holders bargaining alone, the model is silent on whether a coalition of the Netherlands, Japan and South Korea would help; that would require assets that are complements for China and members able to compensate one another.
